\subsection{KRULL 整环的除子类}

设 A 为 Krull 整环。回忆，A 的除子群 D(A) 是由其极值元集合 P(A) 生成的自由交换群（第3节，定理2），且 P(A) 可与 A 的高度为1的素理想集合等同（第6节）；对于 \(p \in P(A)\)，我们把对应于 p 的赋范本性赋值记为 \(v_{p}\)（第4节）；回忆，\(v_{p}\) 的赋值环是 \(A_{p}\)（第4节，命题6后的推论1）。我们把由主除子组成的 D(A) 子群记为 F(A)，把 C(A) = D(A)/F(A) 记为 A 的除子类群（第2节）。

命题14. 设 A 为 Krull 整环，B 为包含 A 的 Krull 整环。假设满足以下条件：

（PDE）对于每个高度为1的素理想 \(\mathfrak{P}\)（属于 \(\mathbf{B}\)），素理想 \(\mathfrak{P} \cap \mathbf{A}\) 为零或具有高度1。

对于 \(\mathfrak{p} \in \mathrm{P}(\mathbf{A})\)，\(\mathfrak{P} \in \mathrm{P}(\mathbf{B})\) 中满足 \(\mathfrak{P} \cap \mathbf{A} = \mathfrak{p}\) 的理想数目有限；记作

\[
i (\mathfrak {p}) = \sum_ {\mathfrak {P} \in \mathrm{P} (\mathrm{B}), \mathfrak {P} \cap \mathrm{A} = \mathfrak {p}} e (\mathfrak {P} / \mathfrak {p}) \mathfrak {P},
\]

其中 \(e(\mathfrak{P} / \mathfrak{p})\) 表示 \(v_{\mathfrak{P}}\) 关于 \(v_{\mathfrak{p}}\) 的分歧指数（第VI章，第8节，第1号）。于是，\(i\) 通过线性性定义了从 D(A) 到 D(B) 的保序同态，并具有以下性质：

\begin{enumerate}
\def\labelenumi{(\alph{enumi})}
\tightlist
\item
  对于每个非零元素 \(x\)（属于 \(\mathbf{A}\) 的分式域），
\end{enumerate}

\[
i \left(\operatorname{div} _ {\mathbf {A}} (x)\right) = \operatorname{div} _ {\mathbf {B}} (x);
\]

\begin{enumerate}
\def\labelenumi{(\alph{enumi})}
\setcounter{enumi}{1}
\tightlist
\item
  对于 D(A) 中的所有 D、D'，
\end{enumerate}

\[
i \left(\sup (\mathbf {D}, \mathbf {D} ^ {\prime})\right) = \sup (i (\mathbf {D}), i (\mathbf {D} ^ {\prime})).
\]

设 \(\mathfrak{p} \in \mathrm{P}(\mathbf{A})\)；取一个非零元素 \(a\)，它属于 \(\mathfrak{p}\)；\(\mathfrak{P} \in \mathrm{P}(\mathbf{B})\) 中包含 \(a\) 的理想数目有限（第6节，定理4）；从而，\(\mathfrak{P} \in \mathrm{P}(\mathbf{B})\) 中满足 \(\mathfrak{P} \cap \mathbf{A} = \mathfrak{p}\) 的理想数目也有限。

现在证明 (a)。由可加性，不妨设 \(x \in \mathbf{A}^* = \mathbf{A} - \{0\}\)。按定义，\(\operatorname{div}_{\mathbf{B}}(x) = \sum_{\mathfrak{P} \in \mathrm{P}(\mathbf{B})} v_{\mathfrak{P}}(x) \cdot \mathfrak{P}\)。对于所有 \(\mathfrak{P} \in \mathrm{P}(\mathbf{B})\)，若 \(v_{\mathfrak{P}}(x) > 0\)，则 \(\mathfrak{P} \cap \mathbf{A}\) 非零（因为 \(x \in \mathfrak{P}\)），由（PDE）可知其具有高度1；令 \(\mathfrak{p} = \mathfrak{P} \cap \mathbf{A}\)，根据分歧指数的定义，\(v_{\mathfrak{P}}(x) = e(\mathfrak{P}/\mathfrak{p})v_{\mathfrak{p}}(x)\)（因为 \(v_{\mathfrak{p}}\) 和 \(v_{\mathfrak{P}}\) 都是赋范的）。又由于 \(\operatorname{div}_{\mathbf{A}}(x) = \sum_{\mathfrak{p} \in \mathrm{P}(\mathbf{A})} v_{\mathfrak{p}}(x) \cdot \mathfrak{p}\)，对于所有满足 \(i(\mathfrak{q}) = 0\) 的 \(\mathfrak{q} \in \mathrm{P}(\mathbf{A})\)，若它不是形如 \(\mathfrak{Q} \cap \mathbf{A}\)（其中 \(\mathfrak{Q} \in \mathrm{P}(\mathbf{B})\)）的理想，即得 (a)。

为证明 (b)，记

\[
\mathrm{D} = \sum_ {\mathfrak {p} \in \mathrm{P} (\mathrm{A})} n (\mathfrak {p}) \cdot \mathfrak {p} \quad \text { and } \quad \mathrm{D} ^ {\prime} = \sum_ {\mathfrak {p} \in \mathrm{P} (\mathrm{A})} n ^ {\prime} (\mathfrak {p}) \cdot \mathfrak {p};
\]

\(\mathfrak{p}\) 在 \(\sup (\mathrm{D},\mathrm{D}^{\prime})\) 中的系数为 \(\sup (n(\mathfrak{p}),n^{\prime}(\mathfrak{p}))\)。设 \(\mathfrak{P}\) 是 \(\mathrm{P(B)}\) 中的元素。若 \(\mathfrak{P}\cap \mathrm{A} = (0)\)，则 \(\mathfrak{P}\) 在 \(i(\mathrm{D})\) 和 \(i(\mathrm{D}^{\prime})\) 中的系数以及在 \(\sup (i(\mathrm{D}),i(\mathrm{D}^{\prime}))\) 中的系数都为零；因此，\(\mathfrak{P}\) 在 \(i(\sup (\mathrm{D},\mathrm{D}^{\prime}))\) 中的系数为零。若 \(\mathfrak{P}\cap \mathrm{A}\neq (0)\)，则它是高度为1的素理想 \(\mathfrak{p}\)（由（PDE））；记 \(e = e(\mathfrak{P} / \mathfrak{p})\)，则 \(\mathfrak{P}\) 在 \(i(\mathrm{D})\)、\(i(\mathrm{D}^{\prime})\) 和 \(i(\sup (\mathrm{D},\mathrm{D}^{\prime}))\) 中的系数分别为 \(en(\mathfrak{p}),en^{\prime}(\mathfrak{p})\) 和 \(e \cdot\sup (n(\mathfrak{p}),n^{\prime}(\mathfrak{p}))\)；而 \(\sup (i(\mathrm{D}),i(\mathrm{D}^{\prime}))\) 中的系数为

\[
\sup (e \cdot n (\mathfrak {p}), e \cdot n ^ {\prime} (\mathfrak {p})) = e \cdot \sup (n (\mathfrak {p}), n ^ {\prime} (\mathfrak {p})).
\]

这证明了 (b)。

在命题14的假设下，由 (a) 可知，i 通过取商定义了一个从 C(A) 到 C(B) 的同态 \(\bar{i}\)，称为典范同态；有时滥用记号，也把它记作 i。

条件（PDE）在以下两种情形中成立：

\begin{enumerate}
\def\labelenumi{(\arabic{enumi})}
\tightlist
\item
  B 整扩于 A；在这种情形，B 的素理想 P 具有高度1，当且仅当 \(p = P \cap A\) 具有高度1。因为 B 中位于 A 的零理想 (0) 之上的唯一素理想是 (0)（第V章，第2节，第1号，命题1的推论1）；若 P 具有高度1，则 \(p \neq 0\)；若 p 不具有高度1，则存在 A 的素理想 \(p'\)，它不同于 (0) 和 p，并且满足
\end{enumerate}

\(0 \subset p' \subset p\)；但由于 B 是整环而 A 是整闭整环，存在 B 的素理想 \(P'\)，使 \(P' \cap A = p'\) 且 \(P' \subset P\)（第V章，第2节，第4号，定理3），这与假设矛盾。反过来，若 p 具有高度1，则不存在不同于 0 和 P 的 B 的素理想 \(P'\)，使 \(0 \subset P' \subset P\)；否则会有 \(0 \subset P' \cap A \subset p\)，并且根据第V章，第2节，第1号，命题1的推论1，\(P' \cap A\) 将不同于 0 和 p。

\begin{enumerate}
\def\labelenumi{(\arabic{enumi})}
\setcounter{enumi}{1}
\tightlist
\item
  B 是 A-平坦模。更确切地说：
\end{enumerate}

命题15. 设 A、B 是 Krull 整环，B 包含 A 且是 A-平坦模。则：

\begin{enumerate}
\def\labelenumi{(\alph{enumi})}
\item
  命题14的条件（PDE）成立；
\item
  对于 \(\mathfrak{a}\)（\(\mathbf{A}\) 的每个除性理想），\(\mathbf{B}\mathfrak{a}\) 是 \(\mathbf{B}\) 中对应于除子 \(i(\mathrm{div}_{\mathbf{A}}(\mathfrak{a}))\) 的除性理想。
\end{enumerate}

为证明 (a)，假设存在 B 的高度为1的素理想 \(\mathfrak{P}\)，使 \(\mathfrak{P} \cap A\) 既非零也不具有高度1。取一个非零元素 \(x \neq 0\)，它属于 \(\mathfrak{P} \cap A\)。A 的高度为1的理想 \(\mathfrak{p}_i\) 中包含 \(x\) 的那些理想数目有限，且没有一个包含 \(\mathfrak{P} \cap A\)；因此存在一个元素 \(y\)，它属于 \(\mathfrak{P} \cap A\)，并且有 \(y \notin \mathfrak{p}_i\)（对所有 \(i\)）（第II章，第2节，第1号，命题2）。于是，\(\operatorname{div}_{A}(x)\) 和 \(\operatorname{div}_{A}(y)\) 是有序群 P(A) 中互素的元素，从而 \(\sup (\operatorname{div}_{A}(x), \operatorname{div}_{A}(y)) = \operatorname{div}_{A}(x) + \operatorname{div}_{A}(y) = \operatorname{div}_{A}(xy)\)；由于理想 \(Ax \cap Ay\) 和 \(Axy\) 都是除性理想，故 \(Ax \cap Ay = Axy\)。由于 B 是 A-平坦模，\(Bx \cap By = Bxy\)（第I章，第2节，第6号，命题6）。这意味着 \(\sup (v_{\mathfrak{P}}(x), v_{\mathfrak{P}}(y)) = v_{\mathfrak{P}}(xy) = v_{\mathfrak{P}}(x) + v_{\mathfrak{P}}(y)\)，这与不等式 \(v_{\mathfrak{P}}(x) > 0\)、\(v_{\mathfrak{P}}(y) > 0\) 矛盾（因为 \(x\) 和 \(y\) 属于 \(\mathfrak{P}\)）。因此 (a) 已由反证法证明。

现在证明 (b)。若 \(a\) 是 \(A\) 的除性理想，则它是两个分式主理想的交（第5节，命题9的推论2），即

\[
a = d ^ {- 1} (\mathrm{A} a \cap \mathrm{A} b)
\]

其中 \(a, b, d\) 属于 \(\mathbf{A}^*\)；由于 \(\mathbf{B}\) 在 A 上平坦，且有 \(\mathbf{A}, \mathbf{B}a = d^{-1}(\mathbf{B}a \cap \mathbf{B}b)\)（第I章，第2节，第6号，命题6），这说明 \(\mathbf{B}a\) 是除性理想。这也说明 \(\mathrm{div}_{\mathbf{B}}(\mathbf{B}a) = \sup (\mathrm{div}_{\mathbf{B}}(a), \mathrm{div}_{\mathbf{B}}(b)) - \mathrm{div}_{\mathbf{B}}(d)\)；利用命题14的 (a) 和 (b)，可见

\[
\begin{array}{r l} \operatorname{div} _ {\mathbf {B}} (\mathrm{Ba}) & = \sup (i (\operatorname{div} _ {\mathbf {A}} (a)), i (\operatorname{div} _ {\mathbf {A}} (b))) - i (\operatorname{div} _ {\mathbf {A}} (d)) \\ & = i (\sup (\operatorname{div} _ {\mathbf {A}} (a), \operatorname{div} _ {\mathbf {A}} (b))) - i (\operatorname{div} _ {\mathbf {A}} (d)) \\ & = i (\operatorname{div} _ {\mathbf {A}} (\mathrm{A} a \cap \mathrm{Ab})) - i (\operatorname{div} _ {\mathbf {A}} (d)) \\ & = i (\operatorname{div} _ {\mathbf {A}} (d ^ {- 1} (\mathrm{A} a \cap \mathrm{Ab}))) = i (\operatorname{div} _ {\mathbf {A}} (\mathrm{a})). \end{array}
\]

推论。设 A 是局部 Krull 整环，B 是离散赋值环，且 B 支配 A 并且是 A-平坦模。则 A 是域或离散赋值环。

设 M 为 B 的极大理想。由（PDE），\(M \cap A\) 为零或具有高度1。由于按假设它是 A 的极大理想，结论由第7节的命题11得到。

评注。在上述两个情形的前一个中，映射 \(i\colon\mathrm{D}(\mathrm{A})\to\mathrm{D}(\mathrm{B})\) 是单射：因为 \(\mathrm{P}(\mathrm{B})\) 的元素构成 \(\mathrm{D}(\mathrm{B})\) 的一组基，且 \(\mathrm{P}(\mathrm{A})\) 中两个不同的理想不可能是 \(\mathrm{A}\) 上某个 \(\mathrm{P}(\mathrm{B})\) 中同一理想的迹，所以只需验证 \(i(\mathfrak{p})\neq0\) 对所有 \(\mathfrak{p}\in\mathrm{P}(\mathrm{A})\) 都成立；现在这由第V章，第2节，第1号，定理1得到。同样可见，若 \(\mathrm{B}\) 是忠实平坦 A-模，则 \(i\) 是单射（第II章，第2节，第5号，命题11的推论4）。

在下文中，我们拟研究某些 Krull 整环有序对 A、B 所对应的从 C(A) 到 C(B) 的典范同态 \(\bar{\imath}\)。

命题16. 设 A 是 Zariski 环，其完备化 \(\hat{\mathbf{A}}\) 是 Krull 整环。则 A 是 Krull 整环，且典范同态 \(\bar{i}\)，从 C(A) 到 C(\(\hat{\mathbf{A}}\))（由于 \(\hat{\mathbf{A}}\) 是 A-平坦模而定义；参见第III章，第3节，第4号，定理3），是单射。

由于 \(\hat{\mathbf{A}}\) 是整环且 \(\mathbf{A} \subset \hat{\mathbf{A}}\)，\(\mathbf{A}\) 是整环。令 \(\mathbf{L}\) 为 \(\hat{\mathbf{A}}\) 的分式域，并令 \(\mathbf{K} \subset \mathbf{L}\) 为 \(\mathbf{A}\) 的分式域。由 \(\mathbf{A} = \hat{\mathbf{A}} \cap \mathbf{K}\)，\(\mathbf{A}\) 是 Krull 整环（第3节，例4）。事实是，\(\bar{i}: \mathbf{C}(\mathbf{A}) \to \mathbf{C}(\hat{\mathbf{A}})\) 是单射，这由命题15(b)以及如下事实得到：若 \(\mathfrak{b}\hat{\mathbf{A}}\) 是主理想，则 \(\mathfrak{b}\) 是主理想（第III章，第3节，第5号，命题9的推论3）。

现在设 A 为 Krull 整环，S 为 A 的一个不含 0 的乘法子集。群 D(A)（分别地，D(S \(^{-1}\) A)）是以 div(p)（分别地，div(S \(^{-1}\) p)）为基的自由交换群，其中 p 遍历 A 的高度为1的素理想集合（分别地，满足 p ∩ S = ∅ 的 A 的高度为1的素理想集合）（第4节，命题6）；并且，若 p ∩ S = ∅，则 i(div(p)) = div(S \(^{-1}\) p)。因此，D(S \(^{-1}\) A) 可识别为 D(A) 的直因子，它由满足 p ∩ S = ∅ 的 div(p) 生成；其补是 D(A) 的自由交换子群，以满足 p ∩ S ≠ ∅ 的 div(p) 为基；我们将这个补记为 G。由于 i: D(A) → D(S \(^{-1}\) A) 是满射，i: C(A) → C(S \(^{-1}\) A) 也是满射；并且：

\[
\mathbf {G} / (\mathbf {G} \cap \mathbf {F} (\mathbf {A})) = (\mathbf {G} + \mathbf {F} (\mathbf {A})) / \mathbf {F} (\mathbf {A}) = \operatorname{Ker} (\bar {i});\tag{5}
\]

因为，若 \(\mathrm{D}(\mathrm{S}^{-1}\mathrm{A})\) 中的一个元素等于 \(\operatorname{div}_{\mathbf{S}^{-1}\mathbf{A}}(x / s)\)，其中 \(x \in \mathbf{A}\) 且 \(s \in \mathbf{S}\)，则它是 \(i\) 作用于主除子 \(\operatorname{div}_{\mathbf{A}}(x)\) 所得的像（命题14）。

现在设 S 由 A 中的元素族 \((p_{i})_{i\in I}\) 生成，且主理想 \(Ap_{i}\) 全都是素理想。于是，若 \(\mathfrak{p}\) 是 A 的高度为1的素理想，且满足 \(\mathfrak{p} \cap \mathrm{S} \neq \varnothing\)，则 \(\mathfrak{p}\) 包含若干个 \(p_{i}\) 的幂的乘积，因而包含其中一个 \(p_{i}\)，记为 \(p_{\alpha}\)；由于 \(\mathrm{Ap}_{\alpha}\) 是非零素理想且 \(\mathfrak{p}\) 具有高度1，得到 \(\mathfrak{p} = \mathrm{Ap}_{\alpha}\)。因此，在上述记号下，\(\mathrm{G} \subset \mathrm{F}(\mathrm{A})\)，由 (5) 可知 \(\bar{i}\) 的核为零。由此证明以下结果：

命题17. 设 A 为 Krull 整环，S 为 A 的一个不含 0 的乘法子集。则典范同态 \(\bar{i}\) 从 C(A) 到 C(S \(^{-1}\) A) 是满射。若进一步 S 由一族元素 \(p_{i}\) 生成，且主理想 Ap \(_{i}\) 全都是素理想，则 \(\bar{i}\) 是双射。

作为公式 (5) 的第二个应用，考虑以下情形：设 R 为 Krull 整环；令 A 为多项式环 A = R{[}X{]}（第9节，命题13），令 S 为 A 的非零常数多项式集合 R - (0)。A 中满足 \(p \cap S \neq \varnothing\) 的高度为1的素理想是形如 \(p_{0}A\) 的那些，其中 \(p_{0}\) 是 R 的高度为1的素理想（第9节，评注）。因此，在上述引入的记号中，通过把 \(\operatorname{div}_{\mathbf{A}}(\mathfrak{p}_{0}\mathbf{A})\) 与 \(\operatorname{div}_{\mathbf{R}}(\mathfrak{p}_{0})\) 对应，G 可识别为 D(R)。另一方面，\(\mathbf{G} \cap \mathbf{F}(\mathbf{A})\) 可识别为 \(\mathbf{F}(\mathbf{R})\)：因为，若 R 中的理想 \(a_{0}\) 在 A 中生成主理想 \(a_{0}\mathbf{A} = f(\mathbf{X})\mathbf{A}\)，其中 A = R{[}X{]}，则 \(f(0) \in a_{0}\mathbf{A}\)，因为 \(a_{0}\mathbf{A}\) 是环 A 的分次理想（按多项式的通常次数分次），从而 \(f(0) \in a_{0}\)；此外，对于 \(a \in a_{0}\)，有 \(a = f(\mathbf{X})g(\mathbf{X})\)，其中 \(g(\mathbf{X}) \in \mathbf{R}\)，比较次数为0的项得 \(a = f(0)g(0)\)；所以 \(a_{0}\) 是 R 中由 \(f(0)\) 生成的主理想。最后，记 R 的分式域为 K，则 \(S^{-1}A\) 可识别为多项式环 K{[}X{]}，它是主理想整环；所以 \(\mathrm{C}(\mathrm{S}^{-1}\mathrm{A}) = (0)\)。因此，由 (5)，\(\mathrm{C}(\mathrm{A}) = \mathrm{Ker}(\bar{i})\) 可识别为 C(R)，从而证明了以下结果：

命题18. 设 R 为 Krull 整环，A 为多项式环 R{[}X{]}。从 C(R) 到 C(R{[}X{]}) 的典范同态是双射。

